In this lecture, we show that the \emph{classical} query complexity of Simon's problem is
$\bigTh{\sqrt{N}}$, where $N = 2^n$, and that the strings measured by Simon's quantum algorithm are
linearly independent with probability at least $\sfrac{1}{4}$.

\section{Simon's Problem, Classically}

In Simon's problem, $f \colon \set{0,1}^n \to \set{0,1}^n$ is given as a black box with the promise
\[f(x) = f(y) \iff x = y \ \text{ or } \ x \oplus y = s,\]
and we have to find $s$ ($f$ is $1$-to-$1$ if $s = 0^n$ and $2$-to-$1$ with period $s$ otherwise).
Let $N = 2^n$. A classical algorithm learns $f$ only through queries $f(q_1), f(q_2), \dots$, and we
may assume that it never queries the same input twice. A \emph{collision} is a pair of queries
$q_i \neq q_j$ with $f(q_i) = f(q_j)$. By the promise a collision gives $s = q_i \oplus q_j$.

\subsection{Deterministic Algorithms}

\Cref{ex:simon-det-0929} gives a deterministic algorithm with $\bigO{\sqrt{N}}$ queries. This is
optimal by \cref{thm:simon-det-lb-1001}.

\subsection{A Randomized Algorithm}

We query random inputs and wait for a collision.

\begin{algorithm}[ht]
    \caption{Randomized classical algorithm for Simon's problem}
    \label{alg:simon-rand-1001}
    \begin{algorithmic}[1]
        \State choose $k$ distinct inputs $q_1, \dots, q_k$ uniformly at random (without replacement) and query them
        \If{$f(q_i) = f(q_j)$ for some $i \neq j$}
        \State \Return $q_i \oplus q_j$
        \Else
        \State \Return $0^n$
        \EndIf
    \end{algorithmic}
\end{algorithm}

\begin{theorem}[Randomized Upper Bound]\label{thm:simon-rand-ub-1001}
    If $f$ is $1$-to-$1$, \cref{alg:simon-rand-1001} is always correct. If $f$ is $2$-to-$1$ with period
    $s$, it is wrong with probability at most $e^{-k(k-1)/(2N)}$. So with
    $k \geq \sqrt{2N \ln N} + 1$ queries (an integer at most $N$ once $n \geq 3$) the algorithm is correct
    with probability at least $1 - \sfrac{1}{N}$. With $k \geq \sqrt{2N} + 1$ queries, which is
    $\bigO{\sqrt{N}}$, this probability is at least $1 - \sfrac{1}{e}$.
\end{theorem}

\begin{proof}
    If $f$ is $1$-to-$1$ there is no collision, and the output $0^n$ is correct. Let $f$ be $2$-to-$1$
    with period $s \neq 0^n$. The algorithm is wrong only if no collision occurs. Call $q_j \oplus s$
    the \emph{mirror} of $q_j$. A collision occurs at step $i$ exactly when $q_i$ is the mirror of an
    earlier query.

    Suppose that no collision occurred among $q_1, \dots, q_{i-1}$. Their $i - 1$ mirrors are distinct,
    and none is among $q_1, \dots, q_{i-1}$, because that would be a collision. Then all of them are
    among the $N - (i-1)$ inputs not yet queried, and $q_i$ is uniform among those. So the
    probability of a collision at step $i$, given none before, is
    $\sfrac{(i-1)}{(N-i+1)} \geq \sfrac{(i-1)}{N}$. Multiplying over the $k$ steps and using
    $1 - x \leq e^{-x}$,
    \begin{align*}
        \prob{\text{no collision among } q_1, \dots, q_k}
         & \leq \prod_{i=1}^{k} \left(1 - \frac{i-1}{N}\right)
        \leq e^{-\sum_{i=1}^{k} (i-1)/N} \\
         & = e^{-k(k-1)/(2N)} \leq e^{-(k-1)^2/(2N)}.
    \end{align*}
    With $k - 1 \geq \sqrt{2N \ln N}$ this is at most $e^{-\ln N} = \sfrac{1}{N}$, and with
    $k - 1 \geq \sqrt{2N}$ it is at most $e^{-1}$.
\end{proof}

\subsection{Lower Bound for Randomized Algorithms via Yao's Principle}

\begin{theorem}[Randomized Lower Bound]\label{thm:simon-rand-lb-1001}
    Let $n \geq 2$. Every randomized classical algorithm that solves Simon's problem with success
    probability at least $\sfrac{2}{3}$ on every input makes more than $\sfrac{\sqrt{N}}{2}$ queries on
    some input. So every such algorithm needs $\bigOm{\sqrt{N}} = \bigOm{\sqrt{2^n}}$ queries.
\end{theorem}

\begin{lemma}[Yao's Minmax Principle, Easy Direction]\label{lem:yao-1001}
    Let $\mu$ be a probability distribution on the inputs. If every deterministic algorithm that makes at
    most $k$ queries errs with probability greater than $\delta$ on an input drawn from $\mu$, then every
    randomized algorithm that makes at most $k$ queries errs with probability greater than $\delta$ on
    some input.
\end{lemma}

\begin{proof}
    A randomized algorithm $R$ is a deterministic algorithm $A_r$ with at most $k$ queries for every
    fixed random string $r$. Let $\mathrm{err}(A, x) \in \set{0, 1}$ indicate that $A$ is wrong on
    input $x$. Then
    \[\max_{x} \prob[r]{\mathrm{err}(A_r, x) = 1}
        \geq \expectation[x \sim \mu]{\prob[r]{\mathrm{err}(A_r, x) = 1}}
        = \expectation[r]{\prob[x \sim \mu]{\mathrm{err}(A_r, x) = 1}}
        > \delta,\]
    since a maximum is at least an average, the two expectations can be exchanged, and every inner
    probability on the right is greater than $\delta$ by assumption.
\end{proof}

\paragraph{The input distribution.} Let $\mathcal{D}$ be the following distribution of $2$-to-$1$
functions. Choose $s$ uniformly among the $N - 1$ nonzero strings. This splits $\set{0,1}^n$ into
$\sfrac{N}{2}$ pairs $\set{x, x \oplus s}$. Assign distinct values to the pairs, uniformly at
random among all such assignments.

\paragraph{Principle of deferred decisions.} We generate the input lazily. Choose $s$. When the
algorithm queries $q_i$: if $q_i = q_j \oplus s$ for an earlier query $q_j$ (a collision), answer
$f(q_j)$. Otherwise answer a uniformly random value among those not used so far ($N - (i-1)$ of them
as long as no collision has occurred). This produces $\mathcal{D}$, because a uniformly random
injective assignment can be built one value at a time, each uniform among the unused ones.

\begin{lemma}[Few Queries Rarely Collide]\label{lem:collision-1001}
    Let $n \geq 2$ and let $A$ be a deterministic algorithm that makes $k \leq \sfrac{\sqrt{N}}{2}$
    queries, all different, on an input drawn from $\mathcal{D}$.
    \begin{enumerate}
        \item The probability that two of the queries form a pair $\set{x, x \oplus s}$ (a collision)
              is at most
              \[\sum_{i=1}^{k} \frac{i-1}{N - 1 - \binom{i-1}{2}} \leq \frac{k(k-1)}{N}.\]
        \item Given that there is no collision among the $k$ queries, and given the $k$ answers, $s$ is
              uniformly distributed over a set of at least $\sfrac{N}{2}$ strings.
    \end{enumerate}
\end{lemma}

\begin{proof}
    Let $Q_i$ be the event that there is no collision among the first $i$ queries, and $C_i$ the
    event $q_i \oplus s \in \set{q_1, \dots, q_{i-1}}$ that the $\nth{i}$ query collides with an
    earlier one.

    \emph{What the algorithm knows about $s$.} Fix $i \leq k + 1$ and condition on $Q_{i-1}$ and on the
    (distinct) answers $f(q_1), \dots, f(q_{i-1})$. Since $A$ is deterministic, the queries
    $q_1, \dots, q_{i-1}$ are the same for every $s$ that leads to these answers. By the lazy
    generation, this sequence of answers has probability $\prod_{j=1}^{i-1} \sfrac{1}{(N - j + 1)}$, the
    same for every $s$ not of the form $q_j \oplus q_l$ with $j < l \leq i - 1$. It is $0$ for an $s$
    of this form ($q_j$ and $q_l$ would have collided). By Bayes' rule, given the answers, $s$ is
    uniform over the nonzero strings that are not among the at most $\binom{i-1}{2}$ strings
    $q_j \oplus q_l$. This is a set of at least $N - 1 - \binom{i-1}{2} \geq N - 1 - \sfrac{k^2}{2} \geq
        N - 1 - \sfrac{N}{8} \geq \sfrac{N}{2}$ candidates (as $i - 1 \leq k \leq \sfrac{\sqrt{N}}{2}$
    and $N \geq 4$). For $i = k + 1$ this is part~2.

    \emph{Part 1.} If $C_i$ is the first collision then $Q_{i-1}$ holds, so by the union bound
    $\prob{\text{collision}} \leq \sum_{i=1}^{k} \prob{C_i \cap Q_{i-1}}$. Given $Q_{i-1}$ and the answers,
    $q_i$ is determined, and $C_i$ occurs exactly when $s$ is one of the $i - 1$ strings
    $q_i \oplus q_1, \dots, q_i \oplus q_{i-1}$. Averaging over the answers,
    \[\prob{C_i \cap Q_{i-1}} \leq \frac{i-1}{N - 1 - \binom{i-1}{2}} \leq \frac{2(i-1)}{N},
        \qquad \text{so} \qquad
        \prob{\text{collision}} \leq \sum_{i=1}^{k} \frac{2(i-1)}{N} = \frac{k(k-1)}{N}. \qedhere\]
\end{proof}

\begin{proof}[Proof of \cref{thm:simon-rand-lb-1001}]
    Let $A$ be a deterministic algorithm with at most $\sfrac{\sqrt{N}}{2}$ queries. We may assume
    that it never repeats a query and makes the same number $k \leq \sfrac{\sqrt{N}}{2}$ of queries on
    every input, because extra queries whose answers are ignored do not change the output. Run it on
    an input drawn from $\mathcal{D}$. It is correct only if its output is $s$. Suppose there is no
    collision. Given the answers, the output is one fixed string, and $s$ is uniform over at least
    $\sfrac{N}{2}$ strings by part~2 of \cref{lem:collision-1001}. So $A$ is correct with probability
    at most $\sfrac{2}{N}$. With part~1,
    \[\prob{A \text{ correct}} \leq \prob{\text{collision}} + \frac{2}{N} \leq \frac{k(k-1) + 2}{N}.\]
    If $k \geq 2$, then $k(k-1) + 2 \leq k^2 \leq \sfrac{N}{4}$ and the bound is at most $\sfrac{1}{4}$.
    If $k \leq 1$, there is no collision and the bound is $\sfrac{2}{N} \leq \sfrac{1}{2}$. So $A$ errs
    with probability at least $\sfrac{1}{2} > \sfrac{1}{3}$ on $\mathcal{D}$. By \cref{lem:yao-1001}
    with $\mu = \mathcal{D}$ and $\delta = \sfrac{1}{3}$, every randomized algorithm with at most
    $\sfrac{\sqrt{N}}{2}$ queries errs with probability greater than $\sfrac{1}{3}$ on some input.
    So an algorithm with success probability at least $\sfrac{2}{3}$ on every input makes more than
    $\sfrac{\sqrt{N}}{2}$ queries on some input.
\end{proof}

\begin{corollary}[Classical Query Complexity of Simon's Problem]\label{thm:simon-det-lb-1001}
    The classical query complexity of Simon's problem is $\bigTh{\sqrt{N}} = \bigTh{\sqrt{2^n}}$, both
    for deterministic algorithms and for randomized algorithms with constant success probability.
    So Simon's quantum algorithm, with $\bigO{n}$ queries, is exponentially faster than every
    classical algorithm.
\end{corollary}

\begin{proof}
    A deterministic algorithm that is always correct is a randomized algorithm with success
    probability $1 \geq \sfrac{2}{3}$, so \cref{thm:simon-rand-lb-1001} gives the lower bound for both
    kinds. The upper bounds are \cref{ex:simon-det-0929} and \cref{thm:simon-rand-ub-1001} (the
    bound $e^{-(k-1)^2/(2N)}$ in its proof is at most $\sfrac{1}{3}$ for $k \geq \sqrt{2N \ln 3} + 1$).
\end{proof}

\begin{corollary}[Decision Version]\label{cor:simon-decision-1001}
    Let $n \geq 2$. Every randomized classical algorithm that decides, with success probability at
    least $\sfrac{2}{3}$ on every input, whether $f$ is $1$-to-$1$ or $2$-to-$1$ makes more than
    $\sfrac{\sqrt{N}}{2}$ queries on some input.
\end{corollary}

\begin{proof}
    Let $\mathcal{P}$ be the uniform distribution on $1$-to-$1$ functions. Let $\mu$ draw $f$ from
    $\mathcal{P}$ or from
    $\mathcal{D}$ with probability $\sfrac{1}{2}$ each. Let $A$ be a deterministic algorithm with
    $k \leq \sfrac{\sqrt{N}}{2}$ queries, all different, as in the proof of
    \cref{thm:simon-rand-lb-1001}. Its answer is a function of the sequence
    $T = (f(q_1), \dots, f(q_k))$ of answers.

    Let $t = (a_1, \dots, a_k)$ have distinct entries. Under $\mathcal{P}$ every answer is uniform
    among the unused values, so $\prob[\mathcal{P}]{T = t} = \prod_{i=1}^{k} \sfrac{1}{(N - i + 1)}$.
    Under $\mathcal{D}$, let $p_i$ be the probability of a collision at step $i$, given no collision
    before and the answers $a_1, \dots, a_{i-1}$. By the proof of \cref{lem:collision-1001},
    $p_i \leq \sfrac{2(i-1)}{N}$. Without a collision the $\nth{i}$ answer is uniform among the
    $N - i + 1$ unused values. By \cref{lem:prodsum-1001} below,
    \begin{align*}
        \prob[\mathcal{D}]{T = t} \geq \prod_{i=1}^{k} \frac{1 - p_i}{N - i + 1}
         & \geq \left(1 - \sum_{i=1}^{k} \frac{2(i-1)}{N}\right) \prob[\mathcal{P}]{T = t} \\
         & = \left(1 - \frac{k(k-1)}{N}\right) \prob[\mathcal{P}]{T = t}
        \geq \frac{3}{4} \prob[\mathcal{P}]{T = t},
    \end{align*}
    because $k(k-1) \leq k^2 \leq \sfrac{N}{4}$.

    Let $\mathcal{T}$ be the set of sequences on which $A$ answers ``$1$-to-$1$'', and let
    $q = \prob[\mathcal{P}]{T \in \mathcal{T}}$. Sequences with a repeated entry have probability $0$
    under $\mathcal{P}$. Summing over $t \in \mathcal{T}$ gives
    $\prob[\mathcal{D}]{T \in \mathcal{T}} \geq \frac{3}{4} q$. So $A$ is correct on an input drawn from
    $\mu$ with probability at most
    \[\frac{1}{2} q + \frac{1}{2}\left(1 - \frac{3}{4} q\right) = \frac{1}{2} + \frac{q}{8} \leq \frac{5}{8} < \frac{2}{3}.\]
    So $A$ errs with probability at least $\sfrac{3}{8} > \sfrac{1}{3}$ on $\mu$. The claim
    follows from \cref{lem:yao-1001}.
\end{proof}

\begin{lemma}\label{lem:prodsum-1001}
    For $a_1, \dots, a_m \in [0, 1]$, $\prod_{i=1}^{m} (1 - a_i) \geq 1 - \sum_{i=1}^{m} a_i$.
\end{lemma}

\begin{proof}
    We use induction on $m$, with equality for $m = 1$. Multiplying the hypothesis for $m - 1$ by
    $1 - a_m \geq 0$ gives
    \[\prod_{i=1}^{m} (1 - a_i) \geq \left(1 - \sum_{i=1}^{m-1} a_i\right)(1 - a_m)
        = 1 - \sum_{i=1}^{m} a_i + a_m \sum_{i=1}^{m-1} a_i \geq 1 - \sum_{i=1}^{m} a_i. \qedhere\]
\end{proof}

\section{Simon's Algorithm Revisited}

In the state \cref{eq:simon-final-0929} just before the measurement, the probability of measuring
$y$ in the control register is the sum of the squared amplitudes of all basis states
$\ket{y}\ket{z}$ with this $y$. This is the squared norm of the vector that multiplies $\ket{y}$:
\[\prob{y} = \norm{\frac{1}{2^n} \sum_{x \in \set{0,1}^n} (-1)^{x \cdot y} \ket{f(x)}}_2^2.\]

\subsection{The \texorpdfstring{$1$}{1}-to-\texorpdfstring{$1$}{1} Case}

\begin{exercisebox}
\begin{exercise}\label{ex:simon-11-1001}
    Suppose that $f$ is $1$-to-$1$. What is the probability of measuring a given string $y$ in Simon's
    circuit, and do these probabilities add up to $1$?
\end{exercise}
\tcblower
\begin{wsolution}
    If $f$ is $1$-to-$1$, the $2^n$ vectors $\ket{f(x)}$ are distinct basis states, so they are
    orthonormal. The vector multiplying $\ket{y}$ then has $2^n$ orthogonal components, each of
    absolute value $\sfrac{1}{2^n}$. Its squared norm is
    \[\prob{y} = \sum_{x \in \set{0,1}^n} \left(\frac{1}{2^n}\right)^2 = 2^n \cdot \frac{1}{2^{2n}} = \frac{1}{2^n}.\]
    So the measured $y$ is uniformly distributed over $\set{0,1}^n$, and the probabilities add up to
    $2^n \cdot \sfrac{1}{2^n} = 1$.
\end{wsolution}
\end{exercisebox}

\subsection{The \texorpdfstring{$2$}{2}-to-\texorpdfstring{$1$}{1} Case}

If $f$ is $2$-to-$1$ with period $s \neq 0^n$, every $y$ with $s \cdot y = 0$ is measured with
probability $\sfrac{1}{2^{n-1}}$ and every $y$ with $s \cdot y = 1$ with probability $0$
(\cref{thm:simon-run-0929}, case 1). There are $2^{n-1}$ strings with $s \cdot y = 0$, so the
probabilities add up to $2^{n-1} \cdot \sfrac{1}{2^{n-1}} = 1$.

\subsection{Solving for \texorpdfstring{$s$}{s}}

\Cref{alg:simon-0929} runs the circuit $n - 1$ times. In the $2$-to-$1$ case the measured strings
$y_1, \dots, y_{n-1}$ give the system of $n - 1$ linear equations
$s \cdot y_1 = 0, \dots, s \cdot y_{n-1} = 0$ over $\set{0,1}$. If the $y_i$ are linearly independent,
the system has rank $n - 1$, and Gaussian elimination gives its only nonzero solution $s'$. The
algorithm then tests $f(0^n) = f(s')$ with two classical queries. If $f$ is $2$-to-$1$, then $s$
solves the system, so $s' = s$ and the test is true. If $f$ is $1$-to-$1$, the test is false because
$s' \neq 0^n$ (\cref{thm:simon-0929}).

\begin{exercisebox}
\begin{exercise}\label{ex:indep-1001}\label{thm:indep-1001}
    Let $V$ be a subspace of $\set{0,1}^n$ of dimension $d \geq 1$, so $\abs{V} = 2^d$, and let
    $y_1, y_2, \dots$ be independent uniformly random elements of $V$. Show that
    \begin{enumerate}
        \item $y_1, \dots, y_d$ are linearly independent with probability
              \[\prod_{i=1}^{d} \left(1 - \frac{1}{2^i}\right)
                  = \left(1 - \frac{1}{2}\right)\left(1 - \frac{1}{2^2}\right)\left(1 - \frac{1}{2^3}\right) \cdots \left(1 - \frac{1}{2^d}\right)
                  \geq \frac{1}{4},\]
        \item $y_1, \dots, y_{d-1}$ are linearly independent with probability
              $\prod_{i=2}^{d} \left(1 - \sfrac{1}{2^i}\right) \geq \sfrac{1}{2}$.
    \end{enumerate}
    Deduce that the strings $y_1, \dots, y_{n-1}$ of Simon's algorithm are linearly independent with
    probability at least $\sfrac{1}{4}$ in the $2$-to-$1$ case
    ($V = \set{y \suchthat s \cdot y = 0}$, $d = n - 1$) and at least $\sfrac{1}{2}$ in the
    $1$-to-$1$ case ($V = \set{0,1}^n$, $d = n$).
\end{exercise}
\tcblower
\begin{proof}
    Suppose that $y_1, \dots, y_j$ are linearly independent, $0 \leq j < d$. Their span consists of the
    $2^j$ vectors $\sum_{i \leq j} c_i y_i$ with $c_i \in \set{0,1}$, all different by independence,
    and it lies inside $V$. The vector $y_{j+1}$ is uniform over the $2^d$ elements of
    $V$ and independent of $y_1, \dots, y_j$. Then $y_1, \dots, y_{j+1}$ are linearly independent exactly
    when $y_{j+1}$ is not in this span, which has probability $1 - \sfrac{2^j}{2^d} = 1 - 2^{-(d-j)}$.
    Multiplying for $j = 0, \dots, m-1$, the first $m$ vectors are linearly independent with
    probability
    \[\prod_{j=0}^{m-1} \left(1 - \frac{1}{2^{d-j}}\right) = \prod_{i=d-m+1}^{d} \left(1 - \frac{1}{2^i}\right),\]
    which gives the two products for $m = d$ and $m = d - 1$. By \cref{lem:prodsum-1001} and the
    geometric series $\sum_{i \geq 2} 2^{-i} = \sfrac{1}{2}$,
    \[\prod_{i=2}^{d} \left(1 - \frac{1}{2^i}\right) \geq 1 - \sum_{i=2}^{d} \frac{1}{2^i} \geq \frac{1}{2},
        \qquad
        \prod_{i=1}^{d} \left(1 - \frac{1}{2^i}\right) = \frac{1}{2} \prod_{i=2}^{d} \left(1 - \frac{1}{2^i}\right) \geq \frac{1}{4}.\]
    In the $2$-to-$1$ case each $y_i$ is uniform over the $2^{n-1}$ strings with $s \cdot y = 0$
    (\cref{thm:simon-run-0929}). They are the solution set of one nontrivial homogeneous linear
    equation, so they form a subspace of dimension $d = n - 1$, and we take $n - 1 = d$ samples. In
    the $1$-to-$1$ case each $y_i$ is uniform over $\set{0,1}^n$ (\cref{ex:simon-11-1001}), $d = n$,
    and we take $d - 1$ samples.
\end{proof}
\end{exercisebox}
